%%O resumo é um texto que deve sumarizar o trabalho.
%%Use frases curtas e autocontidas. 
%%Encadeie-as mostrando o contexto do seu trabalho, o problema, a hipótese ou solução proposta, os pontos importantes do método e dos resultados.
%%O resumo deve ser escrito em um parágrafo único, portanto, é recomendado escrevê-lo após o texto estar pronto. 
%%O resumo é seguido de palavras chaves definidas no arqivo tcc.tex.
O acompanhamento de garantias de produtos é uma demanda recorrente nas relações de consumo, já amplamente atendida por aplicativos voltados ao consumidor final. No entanto, há uma carência de ferramentas dedicadas e de baixa complexidade voltadas para pequenos varejistas e assistências técnicas, que frequentemente recorrem a controles manuais imprecisos ou a sistemas de gestão (ERPs) superdimensionados. Diante deste cenário, este trabalho propõe o desenvolvimento de um sistema web no modelo Micro-SaaS (Software as a Service) focado exclusivamente no controle de garantias. A plataforma adota uma arquitetura multiempresa (multi-tenant), permitindo que diferentes negócios atuem de forma independente no cadastro de clientes, produtos e vendas. O sistema viabiliza o acompanhamento automatizado dos prazos de garantia legal, contratual e estendida, além de estruturar o gerenciamento das solicitações de acionamento de garantia, proporcionando maior organização, centralização de dados e eficiência operacional para o microempreendedor.

\begin{keywords}
Controle de garantias, Micro-SaaS, Arquitetura Multi-Tenant, Sistema Web, Gestão de Varejo.
\end{keywords}